\documentclass[letterpaper, 12pt, titlepage]{article} % Tipo de documento, configurado con: tamaño carta, tamaño de fuente 12pt, inclusión de página de título
\sloppy
\usepackage{wallpaper}
\CenterWallPaper{1.0}{bg.jpg}
% ================================================================
%* PAQUETES BÁSICOS Y CONFIGURACIÓN ESENCIAL
% ================================================================
\usepackage[utf8]{inputenc}         % Codificación de caracteres UTF-8
\usepackage[spanish]{babel}         % Localización en español (guiones, títulos)
\usepackage[left=2cm, right=2cm, top=4.25cm, bottom=3.1cm]{geometry}   % Configuración de márgenes de página
%\usepackage[htt]{hyphenat}

% ================================================================
%* TIPOGRAFÍA Y FORMATO DE TEXTO
% ================================================================
%\usepackage{fontspec}              % Manejo de fuentes con XeLaTeX
%\setmainfont{Myriad Pro} % Fuente principal (Myriad Pro)
\usepackage{courier}
\renewcommand{\familydefault}{\ttdefault}
%\usepackage{helvet}                         % Fuente Sans-Serif (Helvetica)
%\renewcommand{\familydefault}{\sfdefault}   % Establece Helvet como fuente principal
\usepackage{setspace}                       % Control de interlineado (1.5, doble espacio)
\usepackage{parskip}                        % Espaciado entre párrafos (en lugar de sangría)
%\usepackage{ragged2e}                       % Mejor manejo de justificación de texto

\renewcommand{\thesection}{\Roman{section}} % Numeración romana para secciones
\renewcommand{\thesubsection}{\arabic{section}.\arabic{subsection}} % Subsecciones: 1.1, 1.2...
\renewcommand{\thesubsubsection}{\arabic{section}.\arabic{subsection}.\arabic{subsubsection}} % Subsubsecciones: 1.1.1, 1.1.2...

% ================================================================
%* ELEMENTOS GRÁFICOS Y DISEÑO
% ================================================================
\usepackage{graphicx}               % Manejo de imágenes (inserción y escalado)

% ================================================================
%* TABLAS PROFESIONALES
% ================================================================
\usepackage{booktabs}              % Reglas tipográficas para tablas (formato APA)
\usepackage{xltabular}             % Tablas largas que ocupan todo el ancho del texto
\usepackage{array}
\setlength{\extrarowheight}{5pt} % Ajusta este valor según necesites
\usepackage[dvipsnames, table]{xcolor}
% ================================================================
% LISTAS Y ESTRUCTURAS
% ================================================================
\usepackage{enumitem}                   % Personalización avanzada de listas
\setlist{nosep, leftmargin=*, wide=0pt} % Listas ultra compactas

% ================================================================
%* HIPERVÍNCULOS Y METADATOS PDF
% ================================================================
\usepackage[hyphens]{url}
\usepackage{hyperref}             % Hipervínculos y referencias internas
\hypersetup{
	breaklinks=true,                % Configuración de hipervínculos
	colorlinks=true,                % Habilita los enlaces coloreados
	linkcolor=blue,                 % Color de los enlaces internos
	urlcolor=blue,                  % Color de los enlaces externos
	citecolor=blue,                 % Color de las citas
	pdfborder={0 0 0}               % Sin bordes en los enlaces
}

% ================================================================
%* TAREAS PENDIENTES Y ANOTACIONES
% ================================================================
\setlength{\marginparwidth}{2cm} % Soluciona el problema de \todonotes
\usepackage{todonotes} % Paquete principal
\setuptodonotes{
	size=\tiny, % Tamaño de las notas
}
% Comandos personalizados (opcional)
\newcommand{\pendiente}[1]{\todo[color=red!30, inline]{#1}}
\newcommand{\revisar}[1]{\todo[color=yellow!30, inline]{#1}}
\newcommand{\actualizar}[1]{\todo[color=green!30, inline]{#1}}

% ================================================================
%! CONFIGURACIONES FINALES
% ================================================================
\newcommand{\universidad}{URACCAN} % Nombre de la universidad

% ================================================================
% PAQUETE LOREM IPSUM PARA TEXTO DE EJEMPLO, BORRAR EN PRODUCCIÓN
% ================================================================
\usepackage{lipsum} % Generador de texto de ejemplo

% ================================================================
% INICIO DEL DOCUMENTO (contenido principal)
% ================================================================
\begin{document}
	
	\section{DATOS GENERALES}
	
	\begin{xltabular}{\linewidth}{|>\bfseries{l}|X|}
		\hline
		\multicolumn{1}{|c|}{\cellcolor{Cyan} \color{white} \textbf{Ítem}} & \multicolumn{1}{c|}{\cellcolor{Cyan} \color{White} \textbf{Descripción}} \\
		\hline
		\cellcolor{Cyan} \color{White} Nombre de la Institución & \universidad\\
		\hline
		\cellcolor{Cyan} \color{White} Unidad Académica Responsable & Coordinación de Emprendimiento y Economía Creativa \\
		\hline
		\cellcolor{Cyan} \color{White} Nombre del Taller & \\
		\hline
		\cellcolor{Cyan} \color{White} Tipo de Evento & Taller \\
		\hline
		\cellcolor{Cyan} \color{White} Modalidad & Presencial con Apoyo Digital\\
		\hline
		\cellcolor{Cyan} \color{White} Lugar de Realización & URACCAN en su Recinto, CURs y Extensiones \\
		\hline
		\cellcolor{Cyan} \color{White} Fecha(s) de Ejecución & XX de mes de 2026 \\
		\hline
		\cellcolor{Cyan} \color{White} Duración Total & 20 horas \\
		\hline
		\cellcolor{Cyan} \color{White} Número de Participantes & 25 \\
		\hline
		\cellcolor{Cyan} \color{White} Público Meta & Estudiantes, docentes, emprendedores y público en general \\
		\hline
		\cellcolor{Cyan} \color{White} Facilitadores/as & \\
		\hline
		\cellcolor{Cyan} \color{White} Responsable del Evento & Coordinación de Emprendimiento y Economía Creativa \\
		\hline
		\cellcolor{Cyan} \color{White} Requisitos de Ingreso & Conocimientos básicos en X\\
		\hline
		\cellcolor{Cyan} \color{White} Tipo de Certificación & Certificado de Participación \\
		\hline
	\end{xltabular}
	
	\section{Objetivos}

	\subsection{Objetivo General}

	\lipsum[1][1-4]

	\subsection{Objetivos Específicos}
	
	\begin{itemize}
		\item \lipsum[1][1-3]
		\item \lipsum[1][1-3]
		\item \lipsum[1][1-3]
		\item \lipsum[1][1-3]
		\item \lipsum[1][1-3]
	\end{itemize}
	
	\section{RESULTADOS DE APRENDIZAJE ESPERADO}
	
	\begin{itemize}
		\item \lipsum[1][1-4]
		\item \lipsum[1][1-4]
		\item \lipsum[1][1-4]
		\item \lipsum[1][1-4]
		\item \lipsum[1][1-4]
	\end{itemize}
	
	\section{CONTENIDOS}
	
	\begin{xltabular}{\linewidth}{|c|X|c|c|c|}
		\hline
		\textbf{Bloque} & \textbf{Contenidos} & \textbf{HT} & \textbf{HP} & \textbf{TOTAL}\\
		\hline
		Bloque 1 & Tema 1 \begin{itemize}
			\item \lipsum[1][1]
			\item \lipsum[1][1]
			\item \lipsum[1][1]
			\item \lipsum[1][1]
		\end{itemize} & 2 & 2 & 4\\
		\hline
		Bloque 2 & Tema 2 \begin{itemize}
			\item \lipsum[1][1]
			\item \lipsum[1][1]
			\item \lipsum[1][1]
			\item \lipsum[1][1]
		\end{itemize} & 2 & 2 & 4\\
		\hline
		Bloque 3 & Tema 3 \begin{itemize}
			\item \lipsum[1][1]
			\item \lipsum[1][1]
			\item \lipsum[1][1]
			\item \lipsum[1][1]
		\end{itemize} & 2 & 2 & 4\\
		\hline
		Bloque 4 & Tema 4 \begin{itemize}
			\item \lipsum[1][1]
			\item \lipsum[1][1]
			\item \lipsum[1][1]
			\item \lipsum[1][1]
		\end{itemize} & 2 & 2 & 4\\
		\hline
		\textbf{TOTAL} & & & & \\
		\hline
	\end{xltabular}
	
	\section{METODOLOGÍA}
	
	\lipsum[1-3]
	
	\section{CRITERIOS DE EVALUACIÓN}
	
	\lipsum[1-3]
	
	\section{REFERENCIAS}
	
	\hangindent=1in
	Autor, A. (Año). \textit{Título de la obra}. Recuperado de \href{https://holamundo.com}{https://holamundo.com}
\end{document}